\documentclass[11pt]{article}
\usepackage{amsmath, amssymb, geometry, graphicx, xcolor, booktabs, hyperref}
\geometry{letterpaper, margin=1in}

\title{\textbf{Homework 2: First-Order Differential Equations\\(Homogeneous, Bernoulli, and Linear Coefficients)}}
\author{Tashfeen Omran}
\date{September 2026}

\begin{document}

\maketitle
\tableofcontents
\newpage

\section{Part 1: Comprehensive Introduction, Context, and Prerequisites}

\subsection{Core Concepts}
In the study of differential equations, we often encounter first-order equations that are neither separable nor linear in their given form. However, through clever algebraic substitutions, these rigid equations can be transformed into familiar, solvable forms. This chapter focuses on three primary classes of such equations:
\begin{enumerate}
    \item \textbf{Homogeneous Equations:} A differential equation $\frac{dy}{dx} = f(x,y)$ is homogeneous if $f(x,y)$ can be expressed entirely as a function of the ratio $\frac{y}{x}$. In differential form $M(x,y)dx + N(x,y)dy = 0$, it is homogeneous if $M$ and $N$ are homogeneous functions of the same degree (i.e., $M(tx, ty) = t^n M(x,y)$). By substituting $y = vx$ (where $v$ is a function of $x$), the equation is guaranteed to become separable.
    \item \textbf{Bernoulli Equations:} A Bernoulli equation takes the form $\frac{dy}{dx} + P(x)y = Q(x)y^n$. It closely resembles a first-order linear equation, except for the problematic $y^n$ term on the right side. By using the substitution $v = y^{1-n}$, this non-linear equation collapses into a standard linear differential equation, which can then be solved using an Integrating Factor.
    \item \textbf{Equations with Linear Coefficients:} These take the form $\frac{dy}{dx} = G\left(\frac{ax + by + c}{Ax + By + C}\right)$. The presence of the constants $c$ and $C$ prevents the equation from being homogeneous. By executing a translation of axes (shifting the origin) via $x = u + h$ and $y = v + k$, we can eliminate the constants and reduce the equation to a homogeneous one.
\end{enumerate}

\subsection{Intuition and Derivation}
The intuition behind these substitutions lies in the fundamental mathematical strategy of \textit{reducing the unknown to the known}. 
For homogeneous equations, if every term scales equally, the geometry of the slope field is constant along any ray originating from the origin. Therefore, the ratio $v = \frac{y}{x}$ describes the slope of these rays. By tracking how this ratio changes ($v$) instead of how $y$ changes independently, the inherent symmetry of the equation factors out $x$, leaving an equation that can be separated.

For Bernoulli equations, the derivation is a direct consequence of the chain rule. If we divide the entire equation by $y^n$, we get $y^{-n}\frac{dy}{dx} + P(x)y^{1-n} = Q(x)$. Notice that the derivative of $y^{1-n}$ is $(1-n)y^{-n}\frac{dy}{dx}$, which perfectly matches the newly created leading term (up to a constant). Thus, setting $v = y^{1-n}$ completely linearizes the equation.

\subsection{Historical Context and Motivation}
The techniques for solving these specific families of equations were developed in the late 17th and early 18th centuries by mathematicians such as Isaac Newton, Gottfried Wilhelm Leibniz, and the Bernoulli brothers (Jakob and Johann). Johann Bernoulli proposed the "Bernoulli equation" in 1695. The primary motivation was solving physical mechanics problems, such as modeling fluid dynamics, finding orthogonal trajectories of families of curves, and analyzing planetary motion. Mathematicians realized that standard separation of variables failed for many physical phenomena, necessitating the invention of variable transformations to force integrability.

\subsection{Key Formulas}
\begin{itemize}
    \item \textbf{Homogeneous Substitution:} $y = vx \implies \frac{dy}{dx} = v + x\frac{dv}{dx}$.
    \item \textbf{Bernoulli Substitution:} For $\frac{dy}{dx} + P(x)y = Q(x)y^n$, use $v = y^{1-n} \implies \frac{dv}{dx} = (1-n)y^{-n}\frac{dy}{dx}$. The linearized form becomes $\frac{dv}{dx} + (1-n)P(x)v = (1-n)Q(x)$.
    \item \textbf{Linear Combination Substitution:} For $\frac{dy}{dx} = f(ax+by+c)$, use $z = ax+by+c \implies \frac{dz}{dx} = a + b\frac{dy}{dx}$.
\end{itemize}

\subsection{Prerequisites}
To master this topic, you must be fully proficient in:
\begin{itemize}
    \item \textbf{Separation of Variables:} The ability to isolate variables and integrate both sides.
    \item \textbf{Linear Integrating Factors:} Solving $y' + P(x)y = Q(x)$ using $\mu(x) = e^{\int P(x)dx}$.
    \item \textbf{Integration Techniques:} Proficiency in integration by parts, partial fractions, and trigonometric integrals.
    \item \textbf{The Chain Rule:} Deep understanding of implicit differentiation, especially when taking the derivative of a substituted variable with respect to $x$.
\end{itemize}

\newpage
\section{Part 2: Detailed Homework Solutions}

\subsection*{Problem 1}
\textbf{Question:} Identify the equation as homogeneous, Bernoulli, linear coefficients, or of the form $y' = G(ax+by)$.
\[ 6tx \, dx - (x^2 + t^2) \, dt = 0 \]

\textbf{Solution:}
Let us rearrange the equation to express the derivative. Assuming $t$ is the independent variable and $x$ is the dependent variable:
\[ 6tx \frac{dx}{dt} = x^2 + t^2 \implies \frac{dx}{dt} = \frac{x^2 + t^2}{6tx} \]
We test for homogeneity. Let $M(t,x) = 6tx$ and $N(t,x) = -(x^2+t^2)$.
\[ M(\lambda t, \lambda x) = 6(\lambda t)(\lambda x) = \lambda^2 (6tx) = \lambda^2 M(t,x) \]
\[ N(\lambda t, \lambda x) = - ((\lambda x)^2 + (\lambda t)^2) = -\lambda^2(x^2+t^2) = \lambda^2 N(t,x) \]
Since both $M$ and $N$ are homogeneous functions of degree 2, the differential equation is homogeneous.

\textbf{Final Answer:} A. homogeneous

\vspace{0.5cm}
\subsection*{Problem 2}
\textbf{Question:} Identify the equation as homogeneous, Bernoulli, linear coefficients, or of the form $y' = G(ax+by)$.
\[ \cos(3x+5y) \, dy = \sin(3x+5y) \, dx \]

\textbf{Solution:}
Rearrange the equation to isolate $\frac{dy}{dx}$:
\[ \frac{dy}{dx} = \frac{\sin(3x+5y)}{\cos(3x+5y)} = \tan(3x+5y) \]
The right side is a pure function of the linear combination $(3x+5y)$. This perfectly matches the form $y' = G(ax+by)$ where $G(z) = \tan(z)$, $a=3$, and $b=5$.

\textbf{Final Answer:} C. the form $y' = G(ax+by)$

\vspace{0.5cm}
\subsection*{Problem 3}
\textbf{Question:} Use the method for solving homogeneous equations to solve the following differential equation.
\[ (8x^2 - y^2) \, dx + (xy - x^3y^{-1}) \, dy = 0 \]
Ignoring lost solutions, if any, an implicit solution in the form $F(x,y) = C$ is $[\quad] = C$.

\textbf{Solution:}
This is a homogeneous equation. We substitute $y = vx$, which means $dy = v \, dx + x \, dv$.
Substitute $y$ and $dy$ into the original equation:
\[ (8x^2 - (vx)^2) \, dx + \left(x(vx) - x^3(vx)^{-1}\right) (v \, dx + x \, dv) = 0 \]
Simplify the terms inside the parentheses:
\[ (8x^2 - v^2x^2) \, dx + \left(vx^2 - \frac{x^3}{vx}\right) (v \, dx + x \, dv) = 0 \]
\[ x^2(8 - v^2) \, dx + \left(vx^2 - \frac{x^2}{v}\right) (v \, dx + x \, dv) = 0 \]
Divide the entire equation by $x^2$ (assuming $x \neq 0$):
\[ (8 - v^2) \, dx + \left(v - \frac{1}{v}\right) (v \, dx + x \, dv) = 0 \]
Distribute the terms:
\[ (8 - v^2) \, dx + v\left(v - \frac{1}{v}\right) \, dx + x\left(v - \frac{1}{v}\right) \, dv = 0 \]
\[ (8 - v^2) \, dx + (v^2 - 1) \, dx + x\left(v - \frac{1}{v}\right) \, dv = 0 \]
Combine the $dx$ terms:
\[ (8 - v^2 + v^2 - 1) \, dx + x\left(v - \frac{1}{v}\right) \, dv = 0 \]
\[ 7 \, dx + x\left(v - \frac{1}{v}\right) \, dv = 0 \]
Separate the variables $x$ and $v$:
\[ \frac{7}{x} \, dx + \left(v - \frac{1}{v}\right) \, dv = 0 \]
Integrate both sides:
\[ \int \frac{7}{x} \, dx + \int \left(v - \frac{1}{v}\right) \, dv = \int 0 \, dx \]
\[ 7\ln|x| + \frac{v^2}{2} - \ln|v| = C \]
Substitute back $v = \frac{y}{x}$:
\[ 7\ln|x| + \frac{1}{2}\left(\frac{y}{x}\right)^2 - \ln\left|\frac{y}{x}\right| = C \]
Using logarithm properties $\ln\left|\frac{y}{x}\right| = \ln|y| - \ln|x|$:
\[ 7\ln|x| + \frac{y^2}{2x^2} - (\ln|y| - \ln|x|) = C \]
\[ 8\ln|x| - \ln|y| + \frac{y^2}{2x^2} = C \]

\textbf{Final Answer:} $8\ln|x| - \ln|y| + \frac{y^2}{2x^2} = C$

\vspace{0.5cm}
\subsection*{Problem 4}
\textbf{Question:} Use the method for solving homogeneous equations to solve the following differential equation.
\[ (18y^2 - xy) \, dx + x^2 \, dy = 0 \]
Ignoring lost solutions, if any, the general solution is $y = [\quad]$.

\textbf{Solution:}
First, rewrite the equation to isolate $\frac{dy}{dx}$:
\[ x^2 \, dy = (xy - 18y^2) \, dx \implies \frac{dy}{dx} = \frac{xy - 18y^2}{x^2} = \frac{y}{x} - 18\left(\frac{y}{x}\right)^2 \]
This is clearly homogeneous. Substitute $y = vx$, which implies $\frac{dy}{dx} = v + x\frac{dv}{dx}$:
\[ v + x\frac{dv}{dx} = v - 18v^2 \]
Subtract $v$ from both sides:
\[ x\frac{dv}{dx} = -18v^2 \]
Separate the variables:
\[ -\frac{1}{18v^2} \, dv = \frac{1}{x} \, dx \]
Integrate both sides:
\[ -\frac{1}{18} \int v^{-2} \, dv = \int \frac{1}{x} \, dx \]
\[ -\frac{1}{18} \left(-v^{-1}\right) = \ln|x| + C \]
\[ \frac{1}{18v} = \ln|x| + C \]
Substitute back $v = \frac{y}{x}$:
\[ \frac{1}{18(y/x)} = \ln|x| + C \implies \frac{x}{18y} = \ln|x| + C \]
Solve for $y$:
\[ y = \frac{x}{18(\ln|x| + C)} \]

\textbf{Final Answer:} $y = \frac{x}{18(\ln|x| + C)}$

\vspace{0.5cm}
\subsection*{Problem 5}
\textbf{Question:} Use the method for solving homogeneous equations to solve the following differential equation.
\[ \frac{dx}{dt} = \frac{x^2 + t\sqrt{9t^2 + x^2}}{tx} \]
Express the equation in the form $\frac{dx}{dt} = F\left(\frac{x}{t}\right)$.
Ignoring lost solutions, if any, an implicit solution in the form $F(t,x) = C$.

\textbf{Solution:}
First, manipulate the right-hand side to express it strictly in terms of $v = \frac{x}{t}$. Divide the numerator and denominator by $t^2$ (assuming $t > 0$ so $\sqrt{t^2} = t$):
\[ \frac{dx}{dt} = \frac{\frac{x^2}{t^2} + \frac{t}{t^2}\sqrt{9t^2 + x^2}}{\frac{tx}{t^2}} = \frac{\left(\frac{x}{t}\right)^2 + \frac{1}{t}\sqrt{t^2\left(9 + \left(\frac{x}{t}\right)^2\right)}}{\frac{x}{t}} \]
\[ \frac{dx}{dt} = \frac{\left(\frac{x}{t}\right)^2 + \sqrt{9 + \left(\frac{x}{t}\right)^2}}{\frac{x}{t}} = \left(\frac{x}{t}\right) + \frac{\sqrt{9 + \left(\frac{x}{t}\right)^2}}{\frac{x}{t}} \]
So, $F\left(\frac{x}{t}\right) = \frac{x}{t} + \frac{\sqrt{9 + (x/t)^2}}{x/t}$.

Now, substitute $x = vt \implies \frac{dx}{dt} = v + t\frac{dv}{dt}$:
\[ v + t\frac{dv}{dt} = v + \frac{\sqrt{9 + v^2}}{v} \]
Subtract $v$ from both sides:
\[ t\frac{dv}{dt} = \frac{\sqrt{9 + v^2}}{v} \]
Separate the variables:
\[ \frac{v}{\sqrt{9 + v^2}} \, dv = \frac{1}{t} \, dt \]
Integrate both sides. For the left side, use $u$-substitution with $u = 9 + v^2, du = 2v \, dv$:
\[ \frac{1}{2} \int u^{-1/2} \, du = \int \frac{1}{t} \, dt \]
\[ \frac{1}{2} \left(2u^{1/2}\right) = \ln|t| + C \implies \sqrt{9 + v^2} = \ln|t| + C \]
Substitute back $v = \frac{x}{t}$:
\[ \sqrt{9 + \left(\frac{x}{t}\right)^2} = \ln|t| + C \]
To write this as $F(t,x) = C$, we move $\ln|t|$ to the left:
\[ \sqrt{9 + \left(\frac{x}{t}\right)^2} - \ln|t| = C \]

\textbf{Final Answer Box 1:} $F\left(\frac{x}{t}\right) = \frac{x}{t} + \frac{\sqrt{9 + (x/t)^2}}{x/t}$ (or equivalent $\frac{(x/t)^2 + \sqrt{9+(x/t)^2}}{x/t}$) \\
\textbf{Final Answer Box 2:} $\sqrt{9 + \left(\frac{x}{t}\right)^2} - \ln|t| = C$

\vspace{0.5cm}
\subsection*{Problem 6}
\textbf{Question:} Use the method for solving homogeneous equations to solve:
\[ \frac{dy}{d\theta} = \frac{\theta \sec\left(\frac{y}{\theta}\right) + y}{\theta} \]
Ignoring lost solutions, if any, the general solution is $y = [\quad]$.

\textbf{Solution:}
Simplify the right side:
\[ \frac{dy}{d\theta} = \sec\left(\frac{y}{\theta}\right) + \frac{y}{\theta} \]
Let $y = v\theta \implies \frac{dy}{d\theta} = v + \theta\frac{dv}{d\theta}$:
\[ v + \theta\frac{dv}{d\theta} = \sec(v) + v \]
Subtract $v$ from both sides:
\[ \theta\frac{dv}{d\theta} = \sec(v) \]
Separate variables (note that $\frac{1}{\sec(v)} = \cos(v)$):
\[ \cos(v) \, dv = \frac{1}{\theta} \, d\theta \]
Integrate both sides:
\[ \int \cos(v) \, dv = \int \frac{1}{\theta} \, d\theta \]
\[ \sin(v) = \ln|\theta| + C \]
Substitute back $v = \frac{y}{\theta}$:
\[ \sin\left(\frac{y}{\theta}\right) = \ln|\theta| + C \]
Solve for $y$ to find the general explicit solution:
\[ \frac{y}{\theta} = \arcsin(\ln|\theta| + C) \]
\[ y = \theta \arcsin(\ln|\theta| + C) \]

\textbf{Final Answer:} $y = \theta \arcsin(\ln|\theta| + C)$

\vspace{0.5cm}
\subsection*{Problem 7}
\textbf{Question:} Use the method for solving homogeneous equations to solve:
\[ \frac{dy}{dx} = \frac{y(\ln(9y) - \ln(9x) + 2)}{2x} \]
Ignoring lost solutions, an implicit solution in the form $F(x,y) = C$.

\textbf{Solution:}
Use logarithm properties to simplify the numerator: $\ln(9y) - \ln(9x) = \ln\left(\frac{9y}{9x}\right) = \ln\left(\frac{y}{x}\right)$.
\[ \frac{dy}{dx} = \frac{y\left(\ln\left(\frac{y}{x}\right) + 2\right)}{2x} = \frac{1}{2}\left(\frac{y}{x}\right)\left(\ln\left(\frac{y}{x}\right) + 2\right) \]
Let $y = vx \implies \frac{dy}{dx} = v + x\frac{dv}{dx}$:
\[ v + x\frac{dv}{dx} = \frac{1}{2}v(\ln(v) + 2) \]
\[ v + x\frac{dv}{dx} = \frac{1}{2}v\ln(v) + v \]
Subtract $v$ from both sides:
\[ x\frac{dv}{dx} = \frac{1}{2}v\ln(v) \]
Separate variables:
\[ \frac{2}{v\ln(v)} \, dv = \frac{1}{x} \, dx \]
Integrate both sides. For the left side, use $u = \ln(v), du = \frac{1}{v} dv$:
\[ 2 \int \frac{1}{u} \, du = \int \frac{1}{x} \, dx \]
\[ 2\ln|\ln(v)| = \ln|x| + C \]
Substitute back $v = \frac{y}{x}$:
\[ 2\ln\left|\ln\left(\frac{y}{x}\right)\right| = \ln|x| + C \]
\[ 2\ln\left|\ln\left(\frac{y}{x}\right)\right| - \ln|x| = C \]

\textbf{Final Answer:} $2\ln\left|\ln\left(\frac{y}{x}\right)\right| - \ln|x| = C$

\vspace{0.5cm}
\subsection*{Problem 8}
\textbf{Question:} Use the method for solving equations of the form $\frac{dy}{dx} = G(ax+by)$ to solve:
\[ \frac{dy}{dx} = \sqrt{14x + y} - 14 \]
Ignoring lost solutions, an implicit solution in the form $F(x,y) = C$.

\textbf{Solution:}
Let $z = 14x + y$. Differentiating with respect to $x$ yields:
\[ \frac{dz}{dx} = 14 + \frac{dy}{dx} \implies \frac{dy}{dx} = \frac{dz}{dx} - 14 \]
Substitute this into the original differential equation:
\[ \frac{dz}{dx} - 14 = \sqrt{z} - 14 \]
\[ \frac{dz}{dx} = \sqrt{z} \]
Separate variables:
\[ z^{-1/2} \, dz = dx \]
Integrate both sides:
\[ \int z^{-1/2} \, dz = \int 1 \, dx \]
\[ 2z^{1/2} = x + C \]
Substitute back $z = 14x + y$:
\[ 2\sqrt{14x + y} = x + C \]
\[ 2\sqrt{14x + y} - x = C \]

\textbf{Final Answer:} $2\sqrt{14x + y} - x = C$

\vspace{0.5cm}
\subsection*{Problem 9}
\textbf{Question:} Use the method for solving equations of the form $\frac{dy}{dx} = G(ax+by)$ to solve:
\[ \frac{dy}{dx} = 4\sin(20x - 5y) \]
Derive and simplify the differential equation for $\frac{dz}{dx}$ where $z = 20x - 5y$.
Now find the solution to the differential equation.

\textbf{Solution:}
Let $z = 20x - 5y$. Differentiating with respect to $x$ gives:
\[ \frac{dz}{dx} = 20 - 5\frac{dy}{dx} \]
Substitute $\frac{dy}{dx} = 4\sin(z)$ into the derivative:
\[ \frac{dz}{dx} = 20 - 5(4\sin(z)) = 20 - 20\sin(z) = 20(1 - \sin(z)) \]
This fills the first box. Now we separate variables to solve:
\[ \frac{1}{1 - \sin(z)} \, dz = 20 \, dx \]
To integrate the left side, multiply the numerator and denominator by the conjugate $(1 + \sin(z))$:
\[ \int \frac{1 + \sin(z)}{1 - \sin^2(z)} \, dz = \int 20 \, dx \]
\[ \int \frac{1 + \sin(z)}{\cos^2(z)} \, dz = 20x + C \]
Split the integral into two fractions:
\[ \int \left(\frac{1}{\cos^2(z)} + \frac{\sin(z)}{\cos^2(z)}\right) \, dz = 20x + C \]
\[ \int (\sec^2(z) + \sec(z)\tan(z)) \, dz = 20x + C \]
Both are fundamental integration forms:
\[ \tan(z) + \sec(z) = 20x + C \]
Substitute back $z = 20x - 5y$:
\[ \tan(20x - 5y) + \sec(20x - 5y) = 20x + C \]
\[ \tan(20x - 5y) + \sec(20x - 5y) - 20x = C \]

\textbf{Final Answer Box 1:} $\frac{dz}{dx} = 20(1 - \sin(z))$ \\
\textbf{Final Answer Box 2:} $\tan(20x - 5y) + \sec(20x - 5y) - 20x = C$

\vspace{0.5cm}
\subsection*{Problem 10}
\textbf{Question:} Use the method for solving Bernoulli equations to solve:
\[ \frac{dy}{dx} + \frac{y}{x} = 6x^7 y^2 \]
Ignoring lost solutions, find the general solution $y = [\quad]$.

\textbf{Solution:}
This is a Bernoulli equation with $n=2$. Divide the entire equation by $y^2$:
\[ y^{-2}\frac{dy}{dx} + \frac{1}{x}y^{-1} = 6x^7 \]
Let $v = y^{1-2} = y^{-1}$. Then $\frac{dv}{dx} = -y^{-2}\frac{dy}{dx}$, which means $y^{-2}\frac{dy}{dx} = -\frac{dv}{dx}$.
Substitute these into the equation:
\[ -\frac{dv}{dx} + \frac{1}{x}v = 6x^7 \]
Multiply by $-1$ to put it in standard linear form:
\[ \frac{dv}{dx} - \frac{1}{x}v = -6x^7 \]
Find the integrating factor $\mu(x)$:
\[ \mu(x) = e^{\int -\frac{1}{x} dx} = e^{-\ln|x|} = \frac{1}{x} \]
Multiply the standard form equation by $\frac{1}{x}$:
\[ \frac{1}{x}\frac{dv}{dx} - \frac{1}{x^2}v = -6x^6 \]
The left side is the product rule derivative of $(\mu(x) v)$:
\[ \frac{d}{dx}\left(\frac{1}{x}v\right) = -6x^6 \]
Integrate both sides:
\[ \frac{1}{x}v = \int -6x^6 \, dx = -\frac{6}{7}x^7 + C \]
Multiply by $x$:
\[ v = -\frac{6}{7}x^8 + Cx \]
Substitute back $v = y^{-1}$:
\[ \frac{1}{y} = Cx - \frac{6}{7}x^8 \]
Invert to solve for $y$:
\[ y = \frac{1}{Cx - \frac{6}{7}x^8} \]

\textbf{Final Answer:} $y = \frac{1}{Cx - \frac{6}{7}x^8}$

\vspace{0.5cm}
\subsection*{Problem 11}
\textbf{Question:} Use the method for solving Bernoulli equations to solve:
\[ \frac{dy}{dx} + \frac{y}{x-3} = 4(x-3)y^{1/3} \]
Find an expression for $v$ in the form $v = y^{1-n}$. Find the linear diff eq satisfied by $v$. Solve for $v$. General solution $y=$.

\textbf{Solution:}
This is a Bernoulli equation with $n=1/3$.
\textbf{Step 1:} The substitution is $v = y^{1 - 1/3} = y^{2/3}$.
\textbf{Step 2:} Differentiate $v$: $\frac{dv}{dx} = \frac{2}{3}y^{-1/3}\frac{dy}{dx}$.
Divide the original equation by $y^{1/3}$:
\[ y^{-1/3}\frac{dy}{dx} + \frac{1}{x-3}y^{2/3} = 4(x-3) \]
Substitute $y^{-1/3}\frac{dy}{dx} = \frac{3}{2}\frac{dv}{dx}$ and $y^{2/3} = v$:
\[ \frac{3}{2}\frac{dv}{dx} + \frac{1}{x-3}v = 4(x-3) \]
Multiply by $\frac{2}{3}$ to isolate $\frac{dv}{dx}$:
\[ \frac{dv}{dx} + \frac{2}{3(x-3)}v = \frac{8}{3}(x-3) \implies \frac{dv}{dx} = \frac{8}{3}(x-3) - \frac{2}{3(x-3)}v \]
\textbf{Step 3:} Solve for $v$. We use the standard form $\frac{dv}{dx} + \frac{2}{3(x-3)}v = \frac{8}{3}(x-3)$.
Integrating factor:
\[ \mu(x) = e^{\int \frac{2}{3(x-3)} dx} = e^{\frac{2}{3}\ln|x-3|} = (x-3)^{2/3} \]
Multiply the equation by $\mu(x)$:
\[ \frac{d}{dx}\left(v(x-3)^{2/3}\right) = \frac{8}{3}(x-3) \cdot (x-3)^{2/3} = \frac{8}{3}(x-3)^{5/3} \]
Integrate both sides:
\[ v(x-3)^{2/3} = \frac{8}{3} \cdot \frac{3}{8} (x-3)^{8/3} + C = (x-3)^{8/3} + C \]
Solve for $v$:
\[ v = (x-3)^2 + C(x-3)^{-2/3} \]
\textbf{Step 4:} General solution for $y$. Since $v = y^{2/3}$, we have $y = v^{3/2}$:
\[ y = \left((x-3)^2 + C(x-3)^{-2/3}\right)^{3/2} \]

\textbf{Final Answer Box 1:} $v = y^{2/3}$ \\
\textbf{Final Answer Box 2:} $\frac{dv}{dx} = \frac{8}{3}(x-3) - \frac{2}{3(x-3)}v$ \\
\textbf{Final Answer Box 3:} $v = (x-3)^2 + C(x-3)^{-2/3}$ \\
\textbf{Final Answer Box 4:} $y = \left((x-3)^2 + C(x-3)^{-2/3}\right)^{3/2}$

\vspace{0.5cm}
\subsection*{Problem 12}
\textbf{Question:} Use the method for solving Bernoulli equations to solve:
\[ \frac{dy}{dx} + 9y = e^x y^{-8} \]
Ignoring lost solutions, the general solution is $y = [\quad]$.

\textbf{Solution:}
Here $n = -8$. The substitution is $v = y^{1 - (-8)} = y^9$.
Differentiating gives $\frac{dv}{dx} = 9y^8 \frac{dy}{dx}$.
Multiply the entire original equation by $9y^8$:
\[ 9y^8\frac{dy}{dx} + 81y^9 = 9e^x \]
Substitute $v$:
\[ \frac{dv}{dx} + 81v = 9e^x \]
Integrating factor: $\mu(x) = e^{\int 81 dx} = e^{81x}$.
Multiply by $\mu(x)$:
\[ e^{81x}\frac{dv}{dx} + 81e^{81x}v = 9e^{82x} \implies \frac{d}{dx}(v e^{81x}) = 9e^{82x} \]
Integrate both sides:
\[ v e^{81x} = \int 9e^{82x} \, dx = \frac{9}{82}e^{82x} + C \]
Solve for $v$:
\[ v = \frac{9}{82}e^x + Ce^{-81x} \]
Since $v = y^9$:
\[ y^9 = \frac{9}{82}e^x + Ce^{-81x} \implies y = \left(\frac{9}{82}e^x + Ce^{-81x}\right)^{1/9} \]

\textbf{Final Answer:} $y = \left(\frac{9}{82}e^x + Ce^{-81x}\right)^{1/9}$

\vspace{0.5cm}
\subsection*{Problem 13}
\textbf{Question:} Use the method for solving Bernoulli equations to solve:
\[ \frac{dy}{dx} + y^5 x + 5y = 0 \]
Ignoring lost solutions, an implicit solution in the form $F(x,y) = C$.

\textbf{Solution:}
Rearrange into standard Bernoulli form:
\[ \frac{dy}{dx} + 5y = -x y^5 \]
Here $n = 5$. Let $v = y^{1-5} = y^{-4}$.
Differentiate: $\frac{dv}{dx} = -4y^{-5}\frac{dy}{dx}$.
Divide the original arranged equation by $y^5$:
\[ y^{-5}\frac{dy}{dx} + 5y^{-4} = -x \]
Substitute the derivatives and variables:
\[ -\frac{1}{4}\frac{dv}{dx} + 5v = -x \]
Multiply by $-4$:
\[ \frac{dv}{dx} - 20v = 4x \]
Integrating factor: $\mu(x) = e^{\int -20 dx} = e^{-20x}$.
Multiply by $\mu(x)$:
\[ \frac{d}{dx}(v e^{-20x}) = 4xe^{-20x} \]
Integrate both sides using integration by parts on the right. Let $u = 4x, du = 4dx$ and $dw = e^{-20x}dx, w = -\frac{1}{20}e^{-20x}$:
\[ \int 4xe^{-20x} \, dx = -\frac{1}{5}xe^{-20x} - \int -\frac{1}{5}e^{-20x} \, dx = -\frac{1}{5}xe^{-20x} - \frac{1}{100}e^{-20x} + C \]
So,
\[ v e^{-20x} = -\frac{1}{5}xe^{-20x} - \frac{1}{100}e^{-20x} + C \]
Substitute back $v = y^{-4}$:
\[ y^{-4} e^{-20x} = -\frac{1}{5}xe^{-20x} - \frac{1}{100}e^{-20x} + C \]
Move terms to form $F(x,y) = C$:
\[ y^{-4} e^{-20x} + \frac{1}{5}xe^{-20x} + \frac{1}{100}e^{-20x} = C \]

\textbf{Final Answer:} $y^{-4} e^{-20x} + \frac{1}{5}xe^{-20x} + \frac{1}{100}e^{-20x} = C$

\vspace{0.5cm}
\subsection*{Problem 14}
\textbf{Question:} Use the method for solving equations with linear coefficients to solve:
\[ (2x-y)dx + (4x+y-3)dy = 0 \]
Express the related system in terms of $h$ and $k$, choose answer. Find DE $\frac{dv}{du}$ in terms of $v, u$. Find DE $\frac{dz}{du}$. Implicit solution in $u, z$. Implicit in $x, y$.

\textbf{Solution:}
\textbf{Step 1:} Let $x = u+h$ and $y = v+k$. Then $dx = du$ and $dy = dv$. Substitute into the DE:
\[ (2(u+h) - (v+k))du + (4(u+h) + (v+k) - 3)dv = 0 \]
\[ (2u - v + 2h - k)du + (4u + v + 4h + k - 3)dv = 0 \]
Set the constant terms to zero to form a system:
\[ 2h - k = 0 \]
\[ 4h + k - 3 = 0 \implies 4h + k = 3 \]
This matches option A. Solving the system: adding the equations gives $6h = 3 \implies h = 1/2$. Then $k = 2(1/2) = 1$.

\textbf{Step 2:} The DE in $u$ and $v$ simplifies to:
\[ (2u - v)du + (4u + v)dv = 0 \implies (4u+v)dv = -(2u-v)du \implies \frac{dv}{du} = \frac{v-2u}{v+4u} \]

\textbf{Step 3:} Substitute $v = zu \implies \frac{dv}{du} = z + u\frac{dz}{du}$.
\[ z + u\frac{dz}{du} = \frac{zu - 2u}{zu + 4u} = \frac{z-2}{z+4} \]
\[ u\frac{dz}{du} = \frac{z-2}{z+4} - z = \frac{z-2 - z^2 - 4z}{z+4} = \frac{-z^2 - 3z - 2}{z+4} \]
\[ \frac{dz}{du} = \frac{-z^2 - 3z - 2}{u(z+4)} \]

\textbf{Step 4:} Separate and integrate:
\[ \frac{z+4}{z^2+3z+2} dz = -\frac{du}{u} \implies \frac{z+4}{(z+1)(z+2)} dz = -\frac{du}{u} \]
Using partial fractions: $\frac{z+4}{(z+1)(z+2)} = \frac{3}{z+1} - \frac{2}{z+2}$.
\[ \int \left(\frac{3}{z+1} - \frac{2}{z+2}\right) dz = -\int \frac{1}{u} du \]
\[ 3\ln|z+1| - 2\ln|z+2| = -\ln|u| + C_1 \]
\[ \ln\left|\frac{(z+1)^3}{(z+2)^2}\right| + \ln|u| = C_1 \implies \frac{u(z+1)^3}{(z+2)^2} = C \]

\textbf{Step 5:} Back-substitute $z = \frac{v}{u}$:
\[ \frac{u\left(\frac{v}{u}+1\right)^3}{\left(\frac{v}{u}+2\right)^2} = \frac{u \frac{(v+u)^3}{u^3}}{\frac{(v+2u)^2}{u^2}} = \frac{(v+u)^3}{(v+2u)^2} = C \]
Back-substitute $u = x - 1/2$ and $v = y - 1$:
$v+u = (y-1) + (x-1/2) = x + y - 3/2$
$v+2u = (y-1) + 2(x-1/2) = 2x + y - 2$
Thus, $\frac{(x+y-3/2)^3}{(2x+y-2)^2} = C$.

\textbf{Final Answer Box 1:} A. $2h-k=0, 4h+k=3$ \\
\textbf{Final Answer Box 2:} $\frac{dv}{du} = \frac{v-2u}{v+4u}$ \\
\textbf{Final Answer Box 3:} $\frac{dz}{du} = \frac{-z^2-3z-2}{u(z+4)}$ \\
\textbf{Final Answer Box 4:} $\frac{u(z+1)^3}{(z+2)^2} = C$ \\
\textbf{Final Answer Box 5:} $\frac{(x+y-3/2)^3}{(2x+y-2)^2} = C$

\vspace{0.5cm}
\subsection*{Problem 15}
\textbf{Question:} Use the substitution $y = v x^{11}$ to solve the differential equation:
\[ \frac{dy}{dx} = \frac{11y}{x} + \cos\left(\frac{y}{x^{11}}\right) \]
Ignoring lost solutions, an implicit solution in the form $F(x,y) = C$.

\textbf{Solution:}
Differentiate the substitution $y = v x^{11}$ with respect to $x$ using the product rule:
\[ \frac{dy}{dx} = x^{11}\frac{dv}{dx} + 11vx^{10} \]
Substitute $y$ and $\frac{dy}{dx}$ into the original equation:
\[ x^{11}\frac{dv}{dx} + 11vx^{10} = \frac{11(vx^{11})}{x} + \cos(v) \]
\[ x^{11}\frac{dv}{dx} + 11vx^{10} = 11vx^{10} + \cos(v) \]
Cancel out $11vx^{10}$ from both sides:
\[ x^{11}\frac{dv}{dx} = \cos(v) \]
Separate the variables:
\[ \frac{1}{\cos(v)} \, dv = \frac{1}{x^{11}} \, dx \implies \sec(v) \, dv = x^{-11} \, dx \]
Integrate both sides:
\[ \int \sec(v) \, dv = \int x^{-11} \, dx \]
\[ \ln|\sec(v) + \tan(v)| = \frac{x^{-10}}{-10} + C = -\frac{1}{10x^{10}} + C \]
Rearrange to bring variables to one side:
\[ \ln|\sec(v) + \tan(v)| + \frac{1}{10x^{10}} = C \]
Substitute back $v = \frac{y}{x^{11}}$:
\[ \ln\left|\sec\left(\frac{y}{x^{11}}\right) + \tan\left(\frac{y}{x^{11}}\right)\right| + \frac{1}{10x^{10}} = C \]

\textbf{Final Answer:} $\ln\left|\sec\left(\frac{y}{x^{11}}\right) + \tan\left(\frac{y}{x^{11}}\right)\right| + \frac{1}{10x^{10}} = C$

\newpage
\section{Part 3: In-Depth Analysis of Problems and Techniques}

\subsection{A) Problem Types and General Approach}
The homework extensively covers transformations that map non-separable, non-linear equations into solvable forms. The problems fall into distinct categories:

\textbf{1. Homogeneous Equations (Problems 1, 3, 4, 5, 6, 7)}
\textit{Identification:} Every term in the equation shares the same combined degree for $x$ and $y$. For instance, in $x^2 + xy$, both terms are degree 2.
\textit{Strategy:} Substitute $y = vx$. The differential $\frac{dy}{dx}$ becomes $v + x\frac{dv}{dx}$. This substitution factors out the independent variable from the right-hand side, guaranteeing that the equation can be algebraically rearranged into a separable differential equation.

\textbf{2. Equations of the Form $y' = G(ax+by)$ (Problems 2, 8, 9)}
\textit{Identification:} The derivative is a function solely of a linear combination of $x$ and $y$.
\textit{Strategy:} Substitute $z = ax+by$. Taking the derivative yields $\frac{dz}{dx} = a + b\frac{dy}{dx}$. Replacing $\frac{dy}{dx}$ transforms the equation into $\frac{dz}{dx} = a + bG(z)$, which is completely separable.

\textbf{3. Bernoulli Equations (Problems 10, 11, 12, 13)}
\textit{Identification:} Matches the standard linear form but has an extra $y^n$ multiplying the function $Q(x)$ on the right side: $y' + P(x)y = Q(x)y^n$.
\textit{Strategy:} Divide by $y^n$ and substitute $v = y^{1-n}$. The chain rule derivative perfectly matches the new first term, turning the non-linear equation into a linear one solvable with the integrating factor $\mu(x) = e^{\int (1-n)P(x)dx}$.

\textbf{4. Linear Coefficients (Problem 14)}
\textit{Identification:} A rational function of two linear equations $\frac{ax+by+c}{Ax+By+C}$.
\textit{Strategy:} Solve the system of equations $ax+by+c=0$ and $Ax+By+C=0$ to find the intersection $(h,k)$. Shift the origin using $x = u+h$ and $y = v+k$. The constants vanish, leaving a standard homogeneous equation in variables $u$ and $v$.

\textbf{5. Custom Substitutions (Problem 15)}
\textit{Identification:} The problem explicitly provides a non-standard substitution (e.g., $y = vx^{11}$).
\textit{Strategy:} Trust the given substitution. Differentiate it carefully using the product or chain rule, and substitute it into the DE. Vast simplification will inevitably occur.

\subsection{B) Key Algebraic and Calculus Manipulations}
\begin{itemize}
    \item \textbf{Logarithmic Absorption (Problem 3, 7):} Using the property $\ln\left(\frac{y}{x}\right) = \ln|y| - \ln|x|$ to simplify implicit solutions, preventing fractions inside logarithms which makes back-substitution cleaner.
    \item \textbf{Partial Fraction Decomposition (Problem 14):} After reducing the linear coefficient equation to a separable equation in $u$ and $z$, the integral $\int \frac{z+4}{(z+1)(z+2)} dz$ required splitting the denominator into $\frac{A}{z+1} + \frac{B}{z+2}$. This is a crucial trick whenever the denominator is a factorable polynomial.
    \item \textbf{Conjugate Multiplication (Problem 9):} To integrate $\frac{1}{1-\sin(z)}$, multiplying the top and bottom by $1+\sin(z)$ allowed the use of the Pythagorean identity $1-\sin^2(z) = \cos^2(z)$, which split the fraction into standard integrals $\sec^2(z)$ and $\sec(z)\tan(z)$.
    \item \textbf{Integration by Parts with Exponentials (Problem 13):} Solving the linear equation required integrating $4xe^{-20x}$. This demanded integration by parts (LIATE rule: Algebraic $x$ before Exponential $e^{-20x}$).
\end{itemize}

\newpage
\section{Part 4: "Cheatsheet" and Tips for Success}

\subsection{Important Formulas}
\begin{itemize}
    \item \textbf{Homogeneous:} $y = vx \implies y' = v + xv'$
    \item \textbf{Bernoulli:} $y' + P(x)y = Q(x)y^n \implies v = y^{1-n}, v' = (1-n)y^{-n}y'$
    \item \textbf{Linear Coeffs:} $x = u+h, y = v+k$ where $(h,k)$ is the intersection of the two lines.
    \item \textbf{Integrating Factor:} For $v' + P(x)v = Q(x)$, $\mu(x) = e^{\int P(x)dx}$.
\end{itemize}

\subsection{Shortcuts and "Cheats"}
\begin{itemize}
    \item \textbf{The Bernoulli Shortcut:} If you recognize a Bernoulli equation $y' + P(x)y = Q(x)y^n$, you can skip the intermediate division step. The transformed equation will \textit{always} perfectly become $\frac{dv}{dx} + (1-n)P(x)v = (1-n)Q(x)$. 
    \item \textbf{Homogeneous Check:} A quick way to verify if an equation is homogeneous is to replace $x$ with $tx$ and $y$ with $ty$. If the $t$'s cancel out completely from the differential equation, it is homogeneous.
    \item \textbf{Log Rule for Separable DEs:} When you get an equation like $\frac{A}{z} dz = \frac{B}{x} dx$, the solution will always be of the form $z^A = C x^B$. You can skip the intermediate $e^{\ln(...)}$ steps by recognizing this pattern.
\end{itemize}

\subsection{Common Pitfalls}
\begin{itemize}
    \item \textbf{Forgetting the Product Rule:} When substituting $y=vx$, students often incorrectly write $\frac{dy}{dx} = \frac{dv}{dx}$. It must be $\frac{dy}{dx} = v + x\frac{dv}{dx}$ because $v$ is a function of $x$, not a constant!
    \item \textbf{Losing the Constant of Integration Early:} In Bernoulli equations, after finding $v$ using the integrating factor, you must include the $+C$ \textit{before} you back-substitute $v = y^{1-n}$. If you divide by the integrating factor after forgetting $C$, your entire family of solutions is lost.
    \item \textbf{Sign Errors in Integrating Factors:} $\mu(x) = e^{\int -P(x) dx}$ is a common mistake. The formula is strictly $e^{\int +P(x) dx}$. If there is a minus sign in the DE, it belongs to $P(x)$. For example, $y' - \frac{1}{x}y$ means $P(x) = -\frac{1}{x}$.
\end{itemize}

\newpage
\section{Part 5: Conceptual Synthesis and The "Big Picture"}

\subsection{A) Thematic Connections}
The central, unifying theme of this entire chapter is \textbf{Transformation of Variables to Force Linearity or Separability}. In mathematics, when we face a problem we cannot solve, we map it to a space where we \textit{can} solve it, solve it there, and then map the answer back. 
We've seen this theme repeatedly:
\begin{itemize}
    \item In algebra, substituting $u = e^x$ to turn $e^{2x} - 5e^x + 6 = 0$ into a solvable quadratic $u^2 - 5u + 6 = 0$.
    \item In integration, using $u$-substitution or trigonometric substitution to turn a non-integrable function into a standard form.
\end{itemize}
Here, the substitution is elevated. We are not just transforming a variable inside an integral; we are transforming the entire coordinate system of the differential equation. 

\subsection{B) Forward and Backward Links}
\textbf{Backward Links:} This topic is the direct logical culmination of basic integration and the separation of variables. Without a firm grasp of separating $x$ and $y$ onto different sides of an equality, none of the transformations in this chapter yield a usable result.

\textbf{Forward Links:} The concept of changing variables to simplify differential equations lays the groundwork for:
\begin{itemize}
    \item \textbf{Multivariable Calculus (Jacobians):} Transforming regions of integration (like Cartesian to Polar) to simplify double and triple integrals.
    \item \textbf{Higher-Order Differential Equations:} We will soon use substitutions like $y = e^{rt}$ to transform second-order linear equations into simple algebraic characteristic equations.
    \item \textbf{Exact Equations:} Homogeneous equations are actually a subset of exact equations, which introduce the concept of potential functions and gradient fields.
\end{itemize}

\newpage
\section{Part 6: Real-World Application and Modeling}

\subsection{A) Concrete Scenarios in Finance and Economics}
\textbf{1. Macroeconomic Growth (The Solow-Swan Model):} In macroeconomics, modeling the capital accumulation of a country over time is critical. The Solow-Swan model assumes constant returns to scale (a homogeneous production function, often Cobb-Douglas). The resulting differential equation for capital per worker over time is a quintessential \textbf{Bernoulli equation}. Solving this dictates how economies reach a steady-state equilibrium.

\textbf{2. Quantitative Finance (Wealth Dynamics):} If a portfolio is managed such that it is continuously rebalanced between a risky asset and a risk-free asset, the wealth equation can take the form of a linear or Bernoulli equation depending on the utility function of the investor (e.g., in continuous-time Merton portfolio optimization). The transformations learned here are foundational for solving deterministic variants of the Hamilton-Jacobi-Bellman (HJB) equations.

\subsection{B) Model Problem Setup: The Solow-Swan Economic Model}
Let us model the capital per worker $k(t)$ in an economy.
\begin{itemize}
    \item \textbf{Variables:} $t = \text{time}$, $k = \text{capital per worker}$.
    \item \textbf{Parameters:} $s = \text{savings rate}$, $\delta = \text{depreciation rate}$, $\alpha = \text{capital share of output} \ (0 < \alpha < 1)$.
    \item \textbf{Production Function:} The output per worker is modeled by a Cobb-Douglas function $y = k^\alpha$.
\end{itemize}
The fundamental equation of motion for capital is that the change in capital ($\frac{dk}{dt}$) equals investment (savings rate times output) minus depreciation:
\[ \frac{dk}{dt} = s k^\alpha - \delta k \]
Rearranging this yields:
\[ \frac{dk}{dt} + \delta k = s k^\alpha \]
This is a standard \textbf{Bernoulli Differential Equation} with $n = \alpha$. To find the trajectory of the economy's capital, an economist would substitute $v = k^{1-\alpha}$, transforming the complex economic growth model into a solvable linear equation!

\newpage
\section{Part 7: Common Variations and Untested Concepts}

A vital topic within the realm of non-linear first-order substitution equations that was absent from your homework is the \textbf{Riccati Equation}.

\subsection{The Riccati Equation}
A Riccati equation is a non-linear differential equation of the form:
\[ \frac{dy}{dx} = P(x)y^2 + Q(x)y + R(x) \]
It looks like a Bernoulli equation with $n=2$, but it has an extra solitary function of $x$ appended to it ($R(x)$). This equation is notoriously difficult because there is no general elementary substitution that solves it.

\textbf{The Trick:} You can only solve a Riccati equation if you are already given \textit{one specific known solution}, say $y_1(x)$. If you know $y_1(x)$, you can use the substitution:
\[ y = y_1(x) + \frac{1}{v(x)} \]
This substitution will brilliantly cancel out the non-linear $y^2$ terms and reduce the Riccati equation into a standard first-order linear equation in terms of $v$.

\textbf{Example Problem:} Solve $y' = y^2 - \frac{2}{x^2}$ given that $y_1 = \frac{2}{x}$ is a particular solution.
\textbf{Worked Solution:}
Let $y = \frac{2}{x} + \frac{1}{v}$. Taking the derivative:
\[ y' = -\frac{2}{x^2} - \frac{1}{v^2}v' \]
Substitute $y$ and $y'$ into the original Riccati equation:
\[ -\frac{2}{x^2} - \frac{1}{v^2}v' = \left(\frac{2}{x} + \frac{1}{v}\right)^2 - \frac{2}{x^2} \]
Expand the squared term on the right:
\[ -\frac{2}{x^2} - \frac{1}{v^2}v' = \frac{4}{x^2} + \frac{4}{xv} + \frac{1}{v^2} - \frac{2}{x^2} \]
Notice that the $-\frac{2}{x^2}$ terms do not perfectly cancel yet, wait.
Ah, $y_1^2 = \frac{4}{x^2}$. Let's recalculate $\left(\frac{2}{x}\right)^2 - \frac{2}{x^2} = \frac{2}{x^2}$.
Let's check if $y_1 = 2/x$ is actually a solution: $y_1' = -2/x^2$. Right hand side: $(2/x)^2 - 2/x^2 = 4/x^2 - 2/x^2 = 2/x^2$. They do not match! 
Let's use a correct particular solution: $y_1 = \frac{1}{x}$.
Check: $y_1' = -\frac{1}{x^2}$. RHS: $\frac{1}{x^2} - \frac{2}{x^2} = -\frac{1}{x^2}$. Matches!
So let $y = \frac{1}{x} + \frac{1}{v}$. Then $y' = -\frac{1}{x^2} - \frac{v'}{v^2}$.
Substitute:
\[ -\frac{1}{x^2} - \frac{v'}{v^2} = \left(\frac{1}{x} + \frac{1}{v}\right)^2 - \frac{2}{x^2} \]
\[ -\frac{1}{x^2} - \frac{v'}{v^2} = \frac{1}{x^2} + \frac{2}{xv} + \frac{1}{v^2} - \frac{2}{x^2} \]
The $\frac{1}{x^2}$ terms perfectly cancel: $-\frac{1}{x^2} = \frac{1}{x^2} - \frac{2}{x^2}$. We are left with:
\[ -\frac{v'}{v^2} = \frac{2}{xv} + \frac{1}{v^2} \]
Multiply the entire equation by $-v^2$:
\[ v' = -\frac{2}{x}v - 1 \implies v' + \frac{2}{x}v = -1 \]
This is now a basic linear equation! Integrating factor $\mu(x) = x^2$.
\[ (x^2 v)' = -x^2 \implies x^2 v = -\frac{x^3}{3} + C \implies v = -\frac{x}{3} + Cx^{-2} \]
Finally, $y = \frac{1}{x} + \frac{1}{-\frac{x}{3} + Cx^{-2}}$.

\newpage
\section{Part 8: Advanced Diagnostic Testing: "Find the Flaw"}

Below are five problems with complete, step-by-step solutions that contain exactly one critical flaw. Identify the error, explain why it is wrong in one sentence, and provide the correct step and final answer.

\subsection*{Flawed Problem 1 (Homogeneous)}
\textbf{Problem:} Solve $\frac{dy}{dx} = \frac{y^2 + x^2}{xy}$.
\textbf{Flawed Solution:}
1. This is a homogeneous equation. Let $y = vx$.
2. Substitute $y$: $\frac{dy}{dx} = \frac{v^2 x^2 + x^2}{vx^2} = \frac{x^2(v^2+1)}{x^2 v} = \frac{v^2+1}{v}$.
3. Substitute the derivative: $v = \frac{v^2+1}{v}$.
4. Multiply by $v$: $v^2 = v^2 + 1 \implies 0 = 1$.
5. Therefore, there is no solution.

\subsection*{Flawed Problem 2 (Bernoulli)}
\textbf{Problem:} Solve $x\frac{dy}{dx} + y = x^2 y^2$.
\textbf{Flawed Solution:}
1. Divide by $x$: $\frac{dy}{dx} + \frac{1}{x}y = x y^2$. This is Bernoulli with $n=2$.
2. Let $v = y^{-1}$. Differentiating gives $\frac{dv}{dx} = y^{-2}\frac{dy}{dx}$.
3. Divide DE by $y^2$: $y^{-2}\frac{dy}{dx} + \frac{1}{x}y^{-1} = x$.
4. Substitute $v$: $\frac{dv}{dx} + \frac{1}{x}v = x$.
5. Integrating factor $\mu(x) = e^{\int 1/x dx} = x$.
6. $x\frac{dv}{dx} + v = x^2 \implies \frac{d}{dx}(xv) = x^2 \implies xv = \frac{x^3}{3} + C \implies v = \frac{x^2}{3} + \frac{C}{x}$.
7. Final answer: $y = \left(\frac{x^2}{3} + \frac{C}{x}\right)^{-1}$.

\subsection*{Flawed Problem 3 (Integrating Factor)}
\textbf{Problem:} Solve the linear equation $\frac{dv}{dx} - \frac{2}{x}v = 4x^3$.
\textbf{Flawed Solution:}
1. The equation is in standard form.
2. Find the integrating factor: $\mu(x) = e^{\int -\frac{2}{x} dx} = e^{-2\ln(x)}$.
3. Simplify: $e^{-2\ln(x)} = -2x$.
4. Multiply DE by $\mu(x)$: $-2x\frac{dv}{dx} + 4v = -8x^4$.
5. $\frac{d}{dx}(-2xv) = -8x^4 \implies -2xv = -\frac{8}{5}x^5 + C \implies v = \frac{4}{5}x^4 - \frac{C}{2x}$.

\subsection*{Flawed Problem 4 (Linear Combination)}
\textbf{Problem:} Solve $\frac{dy}{dx} = (x+y)^2$.
\textbf{Flawed Solution:}
1. Let $z = x+y$. Then $\frac{dz}{dx} = 1 + \frac{dy}{dx}$.
2. Substitute into DE: $\frac{dz}{dx} = z^2$.
3. Separate variables: $z^{-2} dz = dx$.
4. Integrate: $-z^{-1} = x + C$.
5. Back substitute: $-\frac{1}{x+y} = x + C \implies x+y = -\frac{1}{x+C} \implies y = -\frac{1}{x+C} - x$.

\subsection*{Flawed Problem 5 (Algebraic simplification)}
\textbf{Problem:} Solve the homogeneous equation $\frac{dy}{dx} = \frac{y + \sqrt{x^2+y^2}}{x}$. Assume $x > 0$.
\textbf{Flawed Solution:}
1. Divide by $x$: $\frac{dy}{dx} = \frac{y}{x} + \sqrt{1 + \left(\frac{y}{x}\right)^2}$.
2. Substitute $y = vx \implies v + x\frac{dv}{dx} = v + \sqrt{1+v^2}$.
3. Simplify: $x\frac{dv}{dx} = \sqrt{1+v^2}$.
4. Separate variables: $\frac{dv}{\sqrt{1+v^2}} = \frac{dx}{x}$.
5. Integrate: $\arcsin(v) = \ln(x) + C$.
6. Final answer: $\arcsin\left(\frac{y}{x}\right) = \ln(x) + C$.

\end{document}